% -*- coding: utf-8 -*-
% !TEX program = xelatex
\documentclass[plain,noanswer,medmath]{../../template/jnuexam}
\usepackage{../../template/kysx}

\begin{document}


\examtitle{1996}{5}


\loadproblems{3}{1996P3}%载入试卷三
\loadproblems{4}{1996P4}%载入试卷四


\makepart{填空题}{本题共5小题，每小题3分，满分15分}


\useproblem{4}{1}{1}%【同试卷四第一(1)题】


\useproblem{4}{1}{2}%【同试卷四第一(2)题】


\begin{problem}
$y = \ln(x + \sqrt{1 + x^2}$，则$y''|_{x = \sqrt{3}} =$ \fillin{}．
\end{problem}


\begin{problem}
$5$阶行列式$\begin{vmatrix}
  1 - a &     a &     0 &     0 & 0 \\
     -1 & 1 - a &     a &     0 & 0 \\
      0 &    -1 & 1 - a &     a & 0 \\
      0 &     0 &    -1 & 1 - a & a \\
      0 &     0 &     0 &    -1 & 1 - a
\end{vmatrix} =$ \fillin{}．
\end{problem}


\begin{problem}
一实习生用同一台机器接连独立地制造$3$个同种零件，
第$i$个零件是不合格品的概率$P_i = \frac{1}{i + 1}$（$i = 1,2,3$），
以$X$表示$3$个零件中合格品的个数，则$P\{X = 2\} =$ \fillin{}．
\end{problem}


\makepart{选择题}{本题共5小题，每小题3分，满分15分}


\begin{problem}
设$f'(x_0) = f''(x_0) = 0$，$f'''(x_0) > 0$，则下列选项正确的是\pickout{}
\begin{abcd}
\item $f'(x_0)$是$f'(x)$的极大值．	
\item $f(x_0)$是$f(x)$的极大值．
\item $f(x_0)$是$f(x)$的极小值．
\item $(x_0,f(x_0))$是曲线$y = f(x)$的拐点．
\end{abcd}
\end{problem}


\useproblem{3}{2}{3}%【同试卷三第二(3)题】


\useproblem{4}{2}{3}%【同试卷四第二(3)题】


\useproblem{4}{2}{4}%【同试卷四第二(4)题】


\begin{problem}
设$A$, $B$为任意两个事件，且$A \subset B$，$P(B) > 0$，则下列选项必然成立的是
\begin{abcd}
\item $P(A) < P(A|B)$．
\item $P(A) \le P(A|B)$．
\item $P(A) > P(A|B)$．
\item $P(A) \ge P(A|B)$．
\end{abcd}
\end{problem}


\makepart{}{本题满分6分}%【同试卷四第三题】

\useproblem{4}{3}{0}


\makepart{}{本题满分7分}

\begin{problem}
设$f(x,y) = \int_0^{xy} \e^{- t^2} \dt$，
求$\frac{x}{y}\frac{\pd^2 f}{\pd x^2} - 2\frac{\pd^2 f}{\pdx\pd y}
+ \frac{y}{x}\frac{\pd^2 f}{\pd y^2}$．
\end{problem}


\makepart{}{本题满分6分}%【同试卷四第五题】

\useproblem{4}{5}{0}


\makepart{}{本题满分7分}%【同试卷四第七题】

\useproblem{4}{7}{0}


\makepart{}{本题满分9分}

\begin{problem}
已知一抛物线通过$x$轴上的两点$A(1,0)$，$B(3,0)$．
\begin{enumerate}
\item 求证：两坐标轴与该抛物线所围图形的面积等于$x$轴与该抛物线所围图形的面积；
\item 计算上述两个平面图形绕$x$轴旋转一周所产生的两个旋转体体积之比．
\end{enumerate}
\end{problem}


\makepart{}{本题满分5分}

\begin{problem}
设$f(x)$在$[a,b]$上连续，在$(a,b)$内可导，且$\frac{1}{b - a}\int_a^bf (x)\dx = f(b)$．
求证：在$(a,b)$内至少存在一点$\xi$，使$f'(\xi) = 0$．
\end{problem}


\makepart{}{本题满分9分}

\begin{problem}
已知线性方程组$\begin{cases}
  x_1 + x_2 - 2x_3 + 3x_4 = 0, \\
  2x_1 + x_2 - 6x_3 + 4x_4 = - 1, \\
  3x_1 + 2x_2 + px_3 + 7x_4 = - 1, \\
  x_1 - x_2 - 6x_3 - x_4 = t.
\end{cases}$讨论参数$p$, $t$取何值时，方程组有解？无解？
当有解时，试用其导出组的基础解系表示通解．
\end{problem}


\makepart{}{本题满分7分}

\begin{problem}
设有$4$阶方阵$A$满足条件$|3E + A| = 0$，$AA^T = 2E$，$|A| < 0$，
其中$E$是$4$阶单位阵，求方阵$A$的伴随矩阵$A^*$的一个特征值．
\end{problem}


\makepart{}{本题满分7分}%【同试卷四第十一题】

\useproblem{4}{11}{0}


\makepart{}{本题满分7分}

\begin{problem}
某电路装有三个同种电气元件，其工作状态相互独立，且无故障工作时间都服从参数为$\lambda > 0$的指数分布．
当三个元件都无故障时，电路正常工作，否则整个电路不能正常工作，试求电路正常工作的时间$T$的概率分布．
\end{problem}


\end{document}
